\section{A simplex from a complete-intersection flag}
\label{sec:flag}

We now associate to $(X,L)$ a graded support family with the same dimensions
as the spaces of sections. Its exponent sets lie in a simplex whose volume
is exactly the leading coefficient of the Hilbert polynomial. This equality
of volumes will later ensure that changing the shape of the simplex loses
no asymptotic supply of exponents.

Encoding sections by the least exponents associated with a flag is standard
in the theory of Okounkov bodies; see, for example,
\citet[Lemma~1.3]{LazarsfeldMustata2009}. For our complete-intersection flag,
the required bound follows directly from division by its defining sections
and a degree calculation on the final curve.

For positive real numbers $a_1,\ldots,a_n$, write
\begin{equation}\label{eq:simplex-definition}
 \Delta(a_1,\ldots,a_n)
 =\left\{y\in\R_{\geq0}^n:
              \sum_{i=1}^n\frac{y_i}{a_i}\leq1\right\}.
\end{equation}
Its Euclidean volume is $a_1\cdots a_n/n!$.

\begin{proposition}[The initial simplex]\label{prop:flag-simplex}
Let $X$ be a smooth integral complex projective variety of dimension
$n\geq3$, let $L$ be ample, and put $V=L^n$. Choose a positive integer $D$
such that $DL$ is very ample. There is a graded support family
$S_k\subset\Z_{\geq0}^n$ such that
\begin{align}
 S_k&\subset kP_0,
 &P_0&=\Delta(1/D,\ldots,1/D,D^{n-1}V),
 \label{eq:flag-simplex}\\
 \#S_k&=h^0(X,kL)
       =\frac{V}{n!}k^n+O(k^{n-1}),
 &\vol(P_0)&=\frac{V}{n!}.
 \label{eq:flag-count}
\end{align}
For all integers $k\geq1$, $m\geq0$, and $r\geq1$, if
$\M(S_k)$ separates jets through degree $m$ at some distinct $r$-tuple
in $\Torus^n$, then $H^0(X,kL)$ does so at some distinct $r$-tuple in $X$.
\end{proposition}

\begin{proof}
Choose general sections $\sigma_1,\ldots,\sigma_{n-1}$ of $DL$ so that
their successive zero loci form a smooth integral flag
\[
 X=X_0\supset X_1\supset\cdots\supset X_{n-1}=C,
 \qquad
 X_i=X_{i-1}\cap\{\sigma_i=0\}.
\]
Such a flag exists by Bertini's theorem for very ample hyperplane sections.
Choose $p\in C$ and a local frame $e$ of $L$ near $p$. The functions
$z_i=\sigma_i/e^D$ for $1\leq i<n$ have independent differentials at $p$;
complete them by a function $z_n$ to local analytic coordinates centered
at $p$.

Represent every section of $kL$ by its germ in the frame $e^k$. These
representations are compatible with products in different degrees. They
are injective, since a section on an integral variety vanishing on an open
germ vanishes identically. Order Taylor exponents lexicographically,
starting with the exponent of $z_1$, and let $S_k$ be the set of possible
least exponents of nonzero section germs. Lemma~\ref{lem:least-terms}
gives
\[
 \#S_k=h^0(X,kL),\qquad S_k+S_l\subset S_{k+l}.
\]

We next bound these exponents by global vanishing orders. Let
$\alpha=(\alpha_1,\ldots,\alpha_n)\in S_k$ be the least exponent of a
section $s$. Put $s_0=s$. For $1\leq i<n$, suppose inductively that the residual section on
$X_{i-1}$ has least exponent $(\alpha_i,\ldots,\alpha_n)$. Its local
order along $X_i$ is $\alpha_i$. A restriction to the integral divisor
$X_i$ that vanishes near $p$ vanishes identically. The Cartier divisor
exact sequence therefore allows division by
$\sigma_i|_{X_{i-1}}$, one factor at a time, exactly $\alpha_i$ times.
Define
\[
 s_i=
 \left.
 \frac{s_{i-1}}{(\sigma_i|_{X_{i-1}})^{\alpha_i}}
 \right|_{X_i}
 \in H^0\left(X_i,
       \left(k-D\sum_{j=1}^i\alpha_j\right)L|_{X_i}\right).
\]
The exact local order ensures that $s_i$ is nonzero. Its germ is obtained
by dividing by $z_i^{\alpha_i}$ and setting $z_i=0$, so its least exponent
is $(\alpha_{i+1},\ldots,\alpha_n)$. This completes the induction. At
$i=n-1$ we obtain a nonzero section on $C$ of
\[
 \left(k-D\sum_{i=1}^{n-1}\alpha_i\right)L|_C.
\]
Its vanishing order at $p$ is $\alpha_n$. Since
$\deg(L|_C)=D^{n-1}V$, the degree bounds that order and yields
\[
 \alpha_n\leq
 \left(k-D\sum_{i=1}^{n-1}\alpha_i\right)D^{n-1}V.
\]
Equivalently,
\[
 D\sum_{i=1}^{n-1}\alpha_i+
 \frac{\alpha_n}{D^{n-1}V}\leq k,
\]
which proves \eqref{eq:flag-simplex}. The volume formula and the Hilbert
polynomial asymptotic for an ample line bundle give
\eqref{eq:flag-count}.

It remains to transfer jets from the least monomials to the sections. Fix
$k$. If $S_k$ is empty, the asserted implication has a false premise and
there is nothing to prove. Otherwise, choose a basis with distinct least
exponents. A single positive
weight vector exposes all these least terms, even though the series may
have infinite support. Indeed, put $\lambda_n=1$ and choose the preceding
weights recursively so that
\begin{equation}\label{eq:lex-exposing-weights}
 \lambda_i>
 \max_{\alpha\in S_k}\sum_{j>i}\lambda_j\alpha_j
 \qquad(1\leq i<n).
\end{equation}
If $\beta$ is lexicographically later than $\alpha$ and $i$ is their
first differing coordinate, then
\[
 \lambda\cdot(\beta-\alpha)
 \geq\lambda_i-\sum_{j>i}\lambda_j\alpha_j>0.
\]
Here we used that all Taylor exponents are nonnegative.

Now take a torus tuple at which $\M(S_k)$ separates the specified jets.
The substitution $z_i=s^{\lambda_i}Z_i$ sends this tuple into the chosen
coordinate chart for sufficiently small $s>0$, without identifying any
points. Lemma~\ref{lem:weighted-limit} gives the least monomials as limits
of normalized section germs, and Lemma~\ref{lem:rank-transfer} gives jet
separation by the original sections at the image tuple in $X$.
\end{proof}

The weights in the last paragraph may depend on $k$. In contrast, the flag,
the local frame, and the lexicographic order are fixed once and for all.
Thus Proposition~\ref{prop:flag-simplex} supplies a single graded family,
together with the separate finite-degree interpolation implications needed
to use it.
